\section{RAFT Pipeline}
As explained in the previous chapter, the RAFT pipeline consists of two parts: a RAG pipeline and a fine-tuning component. This section describes how each part of this method is constructed and which techniques are employed for each.
The RAG pipeline is built from several components, such as search, prompt builder, and others.
In this section, each design choice is examined and the rationale behind it is explained.
An overview of the RAG pipeline is presented in Figure~\ref{fig:rag-pipeline}.


\begin{figure}[htbp]
    \centering
    \begin{tikzpicture}[
        node distance=1.4cm,
        every node/.style={
                draw,
                rounded corners=8pt,
                minimum height=0.95cm,
                align=center,
                font=\small\sffamily,
                text width=4.8cm
            },
        arrow/.style={-{Latex}, thick}
        ]
        \node (Input) {User Question};

        \node[below=of Input] (Vector) {Vector Search};
        \node[right=of Input] (Keyword) {Keyword Extractor};
        \node[below=of Keyword] (Lexical) {Lexical Search};

        \node[below= of Vector] (Retrieved) {Retrieved and Scored Content};

        \node[below= of Retrieved] (GatherPrereq) {Gather Prerequisites};
        \node[right=of Retrieved] (GatherSym) {Gather Symbols};

        \node[below= of GatherSym] (GetStatus) {Get User Knowledge Status};

        \node[below=1.6cm of GetStatus] (Build) {Build Prompt};

        % Arrows
        \draw[arrow] (Input) -- (Vector);
        \draw[arrow] (Input) -- (Keyword);
        \draw[arrow] (Keyword) -- (Lexical);

        \draw[arrow] (Vector) -- (Retrieved);
        \draw[arrow] (Lexical) -- (Retrieved);

        \draw[arrow] (Retrieved) -- (GatherPrereq);
        \draw[arrow] (Retrieved) -- (GatherSym);
        \draw[arrow] (Retrieved) -- (Build);

        \draw[arrow] (GatherPrereq) -- (GetStatus);
        \draw[arrow] (GatherSym) -- (GetStatus);
        \draw[arrow] (GetStatus) -- (Build);
    \end{tikzpicture}
    \caption{Overview of the RAG pipeline.}
    \label{fig:rag-pipeline}
\end{figure}
\subsection{Search}
Text search plays a critical role in the RAG pipeline because retrieving relevant content and providing it to the LLM improves result consistency and reduces hallucinations.
In this thesis, two types of search are employed to achieve better results: vector search and lexical search.
Combining these two methods yields a hybrid approach that leverages the strengths of both.

The results of the two search methods are FTML content, which contains many useless HTML tags that cause the LLM's context to become full. Therefore, many of these tags are removed to make the HTML more robust and fit the context size of small models.
\subsubsection{Vector Search}
After constructing the dataset, relevant information needs to be retrieved for the RAG pipeline.
By default, MathHub offers a robust text-based search system; however, it does not support vector-based search.
Therefore, vector search functionality was necessary to enable semantic search over the content of MathHub.
Content was retrieved from the URLs of MathHub and a vector database was built for querying this content.

In this approach, various embedding models were experimented with to observe their impact on query results.
The results indicated that increasing the embedding dimensions beyond 1,000 can be counterproductive, yielding poor results.

\subsubsection{Lexical Search}
MathHub provides a text search function; however, this approach has a significant limitation. It performs pure text search,
meaning it only executes exact string matching rather than identifying keywords from a user query.
Furthermore, users tend to ask questions that describe a concept rather than using specific keywords.

For example, if the user input is ``What does an NP-hard problem mean?'' the pure text search will retrieve
unrelated content with low relevance scores because it requires the exact phrase ``NP-hard'' to appear in the results.

This limitation motivated the development of a keyword extraction mechanism from the user input, on which the text search
is then applied.

\subsection{Gather Prerequisites}

Prerequisites are a critical aspect, as explained in detail in the introduction. Each item retrieved from the search is mapped to its specific prerequisites via the MathHub knowledge base. The
prerequisites and their content are fetched for inclusion in the prompt. The purpose of including prerequisites is to enable the LLM to explain them before addressing the user's question (RQ1).

Additionally, the learner model's dimensions of understanding are fetched for every prerequisite. This is essential because it helps address RQ1.
In the following steps, this is referred to as `prerequisites\_level\_documents`.
The prerequisites are categorized into three sections based on the average value of Bloom's metrics. For beginners, all content should be explained. For medium-level learners, the explanation should be limited to two sentences. For experts, the prerequisites are skipped.
\subsection{Gather Symbols}
Each item retrieved from the search contains associated symbols.
The learner model's dimensions of understanding are fetched for every symbol.
This is a vital point because the model's response should vary based on the user's understanding of the topic.
In the prompt, this is referred to as `symbol\_level`.

\subsection{Building the Prompt}

Constructing the prompt is a critical step, as it must present the data in a structured manner that enables the model to comprehend the relationships between the various components. The prompt structure is as follows:
\begin{quote}
    You are a question-answering model for students. Please provide answers that are easy for students to understand.\\
    Use all material provided in the context.
    I provide the user's status regarding some prerequisites of the topic.
    If the user does not know some prerequisites, cover them before explaining the question to the user.
    Answer in an academic tone.
    User query: \{user query\}\\
    Your response should include the following content:
    \{ retrieved documents\}\\
    The symbols have understanding levels; if the user does not know them, please explain them, otherwise skip them:
    \{symbol\_level\}\\
    This is the prerequisites' status and content; if the user does not know them, please add the content, otherwise skip them:\\
    \{prerequisites\_level\_documents\}
\end{quote}

\section{Fine-tuning}

The RAG pipeline described above is essential because the data varies significantly across courses. Fine-tuning the model is a major component of the RAFT pipeline and distinguishes it from standard RAG.
Fine-tuning enables the model to produce responses with a consistent structure and to better interpret the learner model's structure.
In this thesis, the QLoRA method is employed for fine-tuning.
A dataset comprising questions, corresponding answers, and the learner model status was prepared for this purpose, with the questions sourced from the dataset.
The model selected for fine-tuning is Google Gemma 4 E4B \cite{gemmateam2026gemma4technicalreport}.
This is a small model with only 4 billion parameters.
The model was fine-tuned for one epoch and 20 steps on the fine-tuning dataset (Section~\ref{fine_tune_dataset}).
The model was quantized to 8-bit precision and executed using Ollama with Cot is enabled. The loss log of the fine-tuning process is shown in Figure~\ref{plot:log_loss}.

\begin{figure}[htbp]
    \centering
    \begin{tikzpicture}
        \begin{axis}[
            title={Fine-tuning Loss},
            xlabel={x},
            ylabel={y},
        ]
            \addplot[blue, domain=0:5]{x^2};
        \end{axis}
    \end{tikzpicture}
    \caption{The fine-tuning's loss for 20 steps on the fine-tuning dataset.}
    \label{plot:log_loss}
\end{figure}
